%% ===========================================================================
%%  5. Discussion
%% ===========================================================================

\section{Discussion}
\label{sec:discussion}

Our results demonstrate that CP can function as a user-facing interface layer
between AI music models and their human users. We distill three design
principles and discuss limitations.

\subsection{Design Implications}

\textbf{Prediction set size is an intuitive confidence indicator.} 
Existing approaches to communicating model uncertainty often rely on calibrated confidence displays (e.g., exposing raw probability distributions) or selective prediction (e.g., abstaining entirely when uncertain)~\cite{bhatt2021uncertainty,geifman2017selective}. However, continuous probability distributions are often difficult for lay users to interpret correctly~\cite{padilla2021uncertainty}, and pure selective prediction can disrupt workflows by offering no guidance. CP prediction sets provide a powerful middle ground: set cardinality translates complex model uncertainty into a discrete, actionable format without requiring statistical training. $|\Gamma| = 1$ signals certainty, while $|\Gamma| \ge 4$ signals ambiguity. This aligns with Miller's observation that humans reason more reliably about categorical sets than continuous probabilities~\cite{miller2019explanation}. A UI showing a track's valid prediction set enables users to calibrate trust at a glance.

\textbf{$\alpha$ provides a direct user control.} The significance level maps
to an answerable question: ``What risk of missing the true genre am I willing to
accept?'' This enables a spectrum from conservative mode ($\alpha = 0.01$, 6--7
candidates) for high-stakes curation to exploratory mode ($\alpha = 0.10$, 2--3
candidates) for casual browsing, all within the same system.

\textbf{Distribution-free calibration is essential for predictable interface
behavior.} The Gaussian baseline's mass overcoverage and Bootstrap's marginal
undercoverage under the identical prediction model demonstrate that parametric
calibration assumptions produce unreliable user-facing behavior. Only empirical
$p$-value calibration consistently respects the coverage promise across datasets
and $\alpha$ levels.

\subsection{Limitations and Future Work}

\textbf{Exchangeability requirement.} The coverage guarantee assumes calibration
and test data are exchangeable. When deploying on AI-generated music or novel
distributions, this may be violated. Integrating novelty detection that alerts
users when a track falls outside the calibration distribution is a natural
extension.

\textbf{Static prototypes.} The Linear Probe produces fixed prototypes that
cannot adapt to personal genre ontologies or evolving styles. Personalizing
prototypes via user feedback (genre corrections, listening history) could
improve relevance while preserving the CP guarantee.

\textbf{User studies needed.} The present work establishes algorithmic
foundations. A critical next step is evaluating how prediction sets affect user
trust calibration, decision speed, and task accuracy in realistic music curation
scenarios. We believe the quantitative results---valid coverage, compact sets,
and controllable behavior---provide compelling motivation for such studies.

Beyond music, prediction sets are a general-purpose interface primitive. We
encourage investigation of CP as a user-facing layer in other perceptual domains
where single-label predictions fail to communicate the ambiguity inherent in
human perception.
